% \iffalse meta-comment
%
% hit-flow-deps.dtx 是第三方宏包的集中加载
%
% \fi
%
% \section{第三方依赖}
%
% \changes{v3.2a}{2026/09/14}{终稿与报告的依赖并进一份；报告的列表照 Word 排成段：不加竖距、
%   首行缩进两字、续行回左缘、标号后空一个西文空格}
%
% ^^A ======================================================================
% ^^A Module flow-deps: 第三方宏包集中加载与基础配置
% ^^A ======================================================================
%    \begin{macrocode}
%<*flow-deps>
%    \end{macrocode}
%
% \emph{两档的加载序列逐句对齐，差的地方按 \option{stage} 分岔。} 两边原先各有一份，
% 三十四行里二十七行一字不差；顺序也一致，各自多出来的那几句插在不同位置。
% 所以合成一份之后，共同的那些排在顶层，多出来的各自套一道闸，两边看到的
% \emph{相对}顺序与合并前完全一样。
%
% \emph{顺序有讲究。} \pkg{geometry} 与 \pkg{hyperref} 不在这里：版心要早、
% 超链接要晚，各自单独成模块（\file{hit-flow-geometry.dtx} 与
% \file{hit-flow-hyperlink.dtx}），装配序列上一前一后夹着本文件。
% \pkg{subcaption} 必须排在 \pkg{ccaption} 之前。
%
% 合并之前这两份没并，是因为 \pkg{docstrip} 的输入表：它夹在两串单类文件中间，
% 文件名必须是单类专属的，否则全局表排不出一致的先后。两个类合成一个输出之后
% 这条限制没了。
% \emph{报告的列表先前是 \pkg{enumitem} 的原样。} 终稿那一档早就把列表设成
% 「不加竖距、续行回左缘」了（见 \file{src/flow/hit-final-paragraph.dtx}），可那一整块
% 挂在 \option{stage}\texttt{=final} 下面，报告拿不到：列表项之间多出
% \cs{topsep} 与 \cs{itemsep}，实测深圳本科中期那份的三条列表比原件多占
% 45.83，后面整节跟着下移。
%
% 取值去 \file{docx} 里读的。2026 年 9 月版深圳本科中期原件的那三段（`w:numPr`）：
% 段上\emph{一个} \texttt{w:before}／\texttt{w:after} 都没写，只有
% \texttt{w:firstLineChars="200" w:firstLine="480"}（首行缩进两字）；
% \file{numbering.xml} 里那一级是 \texttt{numFmt=decimal}、
% \texttt{lvlText="\%1."}、\texttt{lvlJc=left}、\texttt{suff=space}——标号左对齐在
% 首行缩进处，后面跟\emph{一个空格}再接正文。
%
% 折成 \pkg{enumitem} 的键：\texttt{leftmargin=0}（续行回左缘）、
% 四个间距归零、\texttt{labelwidth=0pt} 加 \texttt{align=left}（标号左对齐、
% 不占额定列宽）、\texttt{labelsep=0.25em}（Times 的空格，小四下 3.0）、
% \texttt{itemindent} 是「两字加一个 \texttt{labelsep}」——\pkg{enumitem} 把标号
% 摆在 \texttt{itemindent} 减 \texttt{labelsep} 减 \texttt{labelwidth} 处，减出来
% 正好是两字。实测标号左缘 109.05、正文 121.05、续行 85.05，与 iota-hit 对原件的
% 复刻件逐值相同。
%
% \emph{只给正文照终稿排的那一档，别的报告不动。} 读过 \file{docx} 的只有
% 深圳本科那两份；别的报告档（例如深圳博士中期）的「说明」页是类自己用
% \env{enumerate}\oarg{itemsep=-4bp} 写的，那一页原件手上没有，动它没有依据——
% 实测放宽到所有报告会把那一页的续行从悬挂改成回左缘、条距从 29 收到 24。
%
% 终稿那一档还是老值（\texttt{itemindent} 3\,em／3.5\,em，标号左缘 112.07）。
% 学位论文的两份范例里一个列表都没有、指南也没规定，那两个数没有依据，
% 要不要一并改成这一套另说。
%    \begin{macrocode}
\hit@setup@math
\RequirePackage{graphicx}
\RequirePackage{xcolor}
\RequirePackage{tikz}
\RequirePackage{pdfpages}
\includepdfset{fitpaper=true}
\hit@setup@list@packages
\bool_if:NT \g__hit_final_body_bool
  {
    \bool_gset_true:N \g__hit_list_nosep_bool
    \clist_gset:Nn \g__hit_list_itemize_clist
      {
        leftmargin=0em,itemsep=0em,partopsep=0em,parsep=0em,topsep=0em,
        itemindent=\dimexpr2\ccwd+0.25em\relax,
        labelwidth=0pt,labelsep=0.25em,align=left
      }
    \clist_gset:Nn \g__hit_list_enumerate_clist
      {
        leftmargin=0em,itemsep=0em,partopsep=0em,parsep=0em,topsep=0em,
        itemindent=\dimexpr2\ccwd+0.25em\relax,
        labelwidth=0pt,labelsep=0.25em,align=left
      }
  }
\bool_if:NF \g__hit_final_bool { \hit@setup@list@spacing }
\hit@setup@footnote@packages
\sys_if_engine_luatex:TF
  { \RequirePackage[normalem]{ulem} }
  { \RequirePackage{xeCJKfntef} }
\hit@engine@postload:
\RequirePackage{longtable}
\bool_if:NF \g__hit_final_bool
  {
    \hit@setup@table
    \hit@setup@footnote
    \hit@setup@publist
  }
\RequirePackage{booktabs}
\RequirePackage{tabularray}
\hit@setup@table@rules
\bool_if:NT \g__hit_debug_bool
  {
    \RequirePackage { layout }
    \RequirePackage { layouts }
    \RequirePackage { lineno }
  }
\RequirePackage{tabularx}
\RequirePackage{xltabular}
\RequirePackage{flafter}
\RequirePackage{subcaption}%支持子图，必须在 ccaption 之前
\RequirePackage{ccaption} %支持双语标题
\hit@setup@caption
%    \end{macrocode}
%
% \emph{终稿还要子图与索引。} 索引只有学位论文排，中英文各一份。
%    \begin{macrocode}
\bool_if:NT \g__hit_final_bool { \hit@setup@subfigure }
\RequirePackage[makeindex]{splitidx}
\newindex[]{china}
\newindex[]{english}
%    \end{macrocode}
%
% \emph{威海硕士开题的页眉在这里定。} 它得排在页面样式定义之后，页面样式在
% 页眉页脚那个模块里定义，那个模块排在本文件之前。
%    \begin{macrocode}
\bool_if:NF \g__hit_final_bool
  {
    \bool_lazy_all:nT
      {
        { \bool_if_p:N \g__hit_weihai_bool }
        { \bool_if_p:N \g__hit_master_bool }
        { \bool_if_p:N \g__hit_proposal_bool }
      }
      { \pagestyle { hit@weihai@headings@main } }
  }
%</flow-deps>
%    \end{macrocode}
%
% \Finale
